\section{生产落地：怎样让 AI SRE 真正进入值班体系}

\subsection{AI SRE 首先是协作工具，而不是替班工具}

课程里一个很重要的现实判断是：短期内 AI SRE 更像“更强的副驾驶”，而不是“完全替代值班工程师的驾驶员”。

\begin{itemize}
    \item 它最擅长的，是在高压场景下快速整合分散信息
    \item 它能帮助工程师更快看见依赖链、近期变更和潜在 blast radius
    \item 但最终是否回滚、是否扩容、是否宣布重大事故，仍然需要人类判断
\end{itemize}

\begin{importantbox}{值班体系中的关键问题不是 ``能不能自动修''，而是 ``什么时候绝不能自动修''}
在生产环境里，少数错误自动化动作的代价可能远大于人工慢几分钟的代价。因此成熟团队会先定义自动化禁区：哪些服务、哪些时段、哪些动作、哪些权限绝不能让 AI 直接执行。
\end{importantbox}

\subsection{把 AI SRE 接入值班流程的三步法}

\begin{table}[H]
\centering
\begin{tabular}{p{0.18\textwidth} p{0.24\textwidth} p{0.32\textwidth}}
\toprule
阶段 & AI 的角色 & 团队关注点 \\
\midrule
观察阶段 & 汇总告警、日志和 traces & 看它是否真的减少排障时间，而不是制造新噪音 \\
建议阶段 & 提供根因假设和缓解动作建议 & 建立建议的准确率、误报率和可审计性 \\
受控执行阶段 & 自动做低风险、可回滚操作 & 明确权限边界、审批门槛和回退机制 \\
\bottomrule
\end{tabular}
\caption{AI SRE 进入值班体系的典型路径}
\end{table}

\begin{warningbox}{没有高质量运行数据，AI SRE 很容易沦为 ``幻觉放大器''}
如果指标命名混乱、trace 覆盖不全、变更记录缺失、runbook 过时，那么 AI 得到的不是更好的上下文，而是一堆更快被拼接起来的噪音。AI SRE 的效果上限，直接受制于团队原有 observability 基础。
\end{warningbox}

\subsection{本章小结}

AI SRE 要真正进入生产值班体系，必须先从观察和建议做起，再逐步进入受控执行。它的上限由模型能力决定，但它的下限几乎完全由监控基础设施和组织边界决定。

